\chapter{Lecture 1: The Linux Shell and File System}
\begin{center}
\textit{January 4, 2018}
\end{center}

\section{Module 1: The Linux Shell}

\begin{definition}
A \textbf{shell} is an interface to the operating system. The OS talks to the
hardware; the shell lets the user talk to the OS. With a shell we
\begin{enumerate}
  \item manipulate files, and
  \item run programs.
\end{enumerate}
\end{definition}

\subsection{Types of Shells}

\paragraph{Graphical shell.} Windows, macOS and Linux all have one: point and
click, drag and drop. Intuitive to use.

The big disadvantage of a graphical shell is that it is not very powerful:
every feature has to be explicitly implemented by somebody. For example, suppose
a folder contains a lot of files with the extension \code{.c} and we want to
rename all of them to have extension \code{.cc}. Unless someone built a ``bulk
rename'' button, we are stuck doing it one file at a time.

\paragraph{Command-line shell.} We type commands at a \emph{prompt} (basically
a blinking cursor). The biggest disadvantage is the steep learning curve; the big
advantage is that it is very powerful. You do not need someone to have
``supported'' a feature: you combine small commands to build the feature
yourself.

DOS used to be the command-line shell on PCs, but it has gotten stale.

\subsection{History of the Linux Shell}
The Linux command-line shell is a descendant of the original Unix shell (1970s),
called simply \emph{the shell} and written by Stephen Bourne. It had a dozen or
two commands. Over time many shells appeared:
\begin{itemize}
  \item the Korn shell (\code{ksh}),
  \item the C shell (\code{csh}), which became the Turbo C shell (\code{tcsh}).
\end{itemize}
By then the original shell was called the \emph{Bourne shell} (\code{sh}), and
its successor became the \textbf{Bourne Again Shell} (\code{bash}), which is
what we use in CS 246.

To check that you are running bash:
\begin{lstlisting}[style=bash]
$ echo $0
-bash
\end{lstlisting}
(\code{\$0} is the name of the currently running program; the leading \code{-}
means it is a login shell.)

\section{The Linux File System}

The file system is similar to the one on Windows. In Linux, there are files
which can contain other files.

\begin{definition}
A \textbf{directory} is a file that can contain other files. Other than that,
Linux does not really distinguish between files and directories.
\end{definition}

The Linux file system is a \textbf{tree}. The root of the tree is a directory
called the \textbf{root directory}, written \code{/}.

\begin{center}
\begin{tikzpicture}[level distance=11mm,
  level 1/.style={sibling distance=30mm},
  level 2/.style={sibling distance=18mm},
  every node/.style={font=\ttfamily\small}]
\node {/}
  child {node {bin}}
  child {node {usr}
    child {node {share}
      child {node {dict}
        child {node {words}}}}
    child {node {bin}}}
  child {node {u}
    child {node {nanaeem}}};
\end{tikzpicture}
\end{center}

\begin{definition}
A \textbf{path} specifies the location of a file in the file system. An
\textbf{absolute path} begins at the root directory, e.g.\
\code{/usr/share/dict/words}. A \textbf{relative path} starts at the current
directory.
\end{definition}

If \code{pwd} is \code{/Users/nanaeem}, then the absolute path
\code{/Users/nanaeem/cs246/a1} names the same directory as the relative path
\code{cs246/a1}.

\begin{aside}
Special directory names:
\begin{itemize}
  \item \code{.} is the current directory,
  \item \code{..} is the parent directory, \code{../..} the grandparent,
  \item \code{\textasciitilde} is your home directory, and
        \code{\textasciitilde userid} is \code{userid}'s home directory.
\end{itemize}
\end{aside}

\section{Basic Commands}

\begin{center}
\begin{tabular}{@{}lp{0.62\textwidth}@{}}
\toprule
Command & Meaning \\
\midrule
\code{clear} & clear the screen \\
\code{ls} & list the (non-hidden) files in the current directory \\
\code{ls -a} & list all files, including hidden ones (names starting with \code{.}) \\
\code{ls -l} & long listing (permissions, owner, size, date) \\
\code{pwd} & print the present working directory \\
\code{cd <path>} & change directory to \code{<path>} \\
\code{cd} & with no argument: go to your home directory \\
\code{cd -} & go back to the previous directory \\
\bottomrule
\end{tabular}
\end{center}

The \code{-a} in \code{ls -a} is called an \emph{option} (or \emph{flag}); it
modifies the behaviour of the command. Options can usually be combined, e.g.\
\code{ls -la}.

\begin{lstlisting}[style=bash]
$ cd /Users/nanaeem/cs246/a1    # absolute path
$ cd ..                          # now in /Users/nanaeem/cs246
$ cd a1                          # relative path, back in a1
$ pwd
/Users/nanaeem/cs246/a1
\end{lstlisting}

\section*{Summary}
\begin{itemize}
  \item A shell is an interface to the OS; command-line shells trade a learning
        curve for power. We use \code{bash}.
  \item The file system is a tree rooted at \code{/}; directories are files that
        contain files.
  \item Absolute paths start at \code{/}; relative paths start at the current
        directory, with \code{.}, \code{..}, \code{\textasciitilde} as shortcuts.
  \item First commands: \code{clear}, \code{ls}, \code{pwd}, \code{cd}.
\end{itemize}
